\sectionguard
\section{The Silence Problem and the Controversy of 1556}

\subsection{What the problem is, stated exactly}

Between December 1531 and the printing of 1648 lies a hundred and
seventeen years in which the apparition narrative does not appear in
any document that the modern critical literature has been able to
authenticate as of that period. The problem is not that the record of
early New Spain is thin. It is unusually thick: episcopal
correspondence, mendicant chronicles, Nahuatl annals, testaments,
viceregal letters, and the great ethnographic project of Bernardino
de Sahag\'un all survive in quantity, and many of these documents
describe precisely the sort of event the narrative reports---heavenly
favours to Indian converts, wonders at shrines, the fortunes of the
Tepeyac hermitage itself.

Two further features sharpen the difficulty. First, the persons the
narrative makes central are among the best-documented figures of the
period: Juan de Zum\'arraga, first bishop and archbishop of Mexico,
left a large body of writings. Second, the one occasion on which the
Tepeyac cult is documented as a subject of formal ecclesiastical
controversy---September 1556---produced a file of sworn testimony in
which, on the reading of the critical historians, nobody mentions the
apparition at all.

The argument built on these facts is an argument from silence, and
its proponents have always known it. Its classical formulation is
Joaqu\'in Garc\'ia Icazbalceta's, in the letter he wrote in 1883 at
the express and repeated command of his archbishop, Pelagio Antonio
Labastida y D\'avalos, and asked him not to circulate. Because the
argument is so often reported second-hand and so rarely quoted, and
because its author was a devout Catholic layman, the founding
historian of Mexican colonial bibliography, and manifestly reluctant
to write at all, it is worth having in his own words.

\begin{witnesstext}{Joaqu\'in Garc\'ia Icazbalceta, \work{Carta
acerca del origen de la imagen de Nuestra Se\~nora de Guadalupe de
M\'exico}, to Archbishop Labastida, October 1883, \S\S9--10 and 12
(Spanish, as edited)}
9.---Ca\'imos ya en el \term{argumento negativo}, tan impugnado por
los apologistas de la Aparici\'on, sin duda porque conocen que no
puede haber otro contra un hecho \term{que no pas\'o}\ldots\par
10.---La fuerza del argumento negativo consiste principalmente en que
el silencio sea universal, y que los autores alegados hayan escrito
de asuntos que ped\'ian una menci\'on del suceso que callaron. Ambas
circunstancias concurren en los documentos anteriores al padre
S\'anchez\ldots\par
12.---El primer testigo de la Aparici\'on debiera ser el
ilustr\'isimo se\~nor Zum\'arraga\ldots\ Pero en los muchos escritos
suyos que conocemos no hay la m\'as ligera alusi\'on al hecho o a
las ermitas: ni siquiera se encuentra una sola vez el nombre de
Guadalupe. Tenemos sus libros de doctrina, cartas, pareceres, una
exhortaci\'on pastoral, dos testamentos y una informaci\'on acerca de
sus buenas obras\ldots\ Pero nada, absolutamente nada en parte
alguna.\endnote{Joaqu\'in Garc\'ia Icazbalceta, \work{Carta acerca
del origen de la imagen de Nuestra Se\~nora de Guadalupe de
M\'exico}, \S\S9--10, 12. Quoted 2026-07-25 from the Biblioteca
Virtual Miguel de Cervantes edition,
\url{https://www.cervantesvirtual.com/obra-visor/carta-acerca-del-origen-de-la-imagen-de-nuestra-senora-de-guadalupe-de-mexico--0/html/b28296be-844c-4ebe-a169-2732231ce572_2.html},
and collated the same day against the first printed edition of 1896
in the Harvard copy digitized by Google at
\url{https://archive.org/details/cartaacercadelo00icazgoog} (Internet
Archive metadata: \texttt{possible-copyright-status:
NOT\_IN\_COPYRIGHT}). The readings agree; the Cervantes Virtual
transcription silently expands the 1896 imprint's abbreviations
(``Ilmo.\ Sr.\ Zum\'arraga,'' ``V.~S.~I.,'' ``P.~S\'anchez''), and
it is that expanded transcription which is quoted here, as this note
records. The paragraph numbering is the letter's own. The Cervantes
Virtual running head prints ``(Octubre 1833)''; the date is 1883, as
its own preface and the documented history of the letter require, and
the figure is treated here as a transcription error rather than
silently corrected inside the quotation. Public domain (Garc\'ia
Icazbalceta d.~1894). Ellipses mark omissions; nothing internal is
altered.}
\end{witnesstext}

Garc\'ia Icazbalceta then extends the inventory: Archbishop
Alonso de Mont\'ufar's own \work{Descripci\'on del Arzobispado} of
1569--1570, which lists the churches of the city and does not mention
the Guadalupe hermitage; Motolin\'ia's \work{Historia de los indios}
of 1541, which recounts heavenly favours to Indians and never names
Guadalupe; the letter of Bishop Garc\'es to Paul~III; the
correspondence of Pedro de Gante, Ram\'irez de Fuenleal and Viceroy
Mendoza; Bartolom\'e de las Casas, who knew Zum\'arraga personally
and to whom the story would have been a gift in his defence of the
Indians; and Jer\'onimo de Mendieta, who wrote three chapters on
Zum\'arraga's life, reported a 1576 apparition to an Indian of
Xochimilco which the visionary told him himself, and said nothing of
Tepeyac.\endnote{Garc\'ia Icazbalceta, \work{Carta}, \S\S13--16, same
witness. The chain of silences is his; this study reports it as his
argument and has not independently examined Motolin\'ia, Mendieta,
Las Casas, the Garc\'es letter, or the 1569--1570 \work{Descripci\'on}
(which Garc\'ia Icazbalceta says he owned in the original).}

To this is added the passage that both schools have always had to
confront: Sahag\'un's note, in the appendix on surviving
superstitions in Book~XI of the \work{Historia general}, that at the
hill called Tepeacac there had formerly been a temple of a goddess
called Tonantzin, ``our mother,'' and that the Indians now came to
the church of Our Lady there still saying that they went to
Tonantzin---which he regarded as a satanic device to keep idolatry
alive under a new name. A missionary who thought the Tepeyac
devotion needed that warning is a difficult witness for a tradition
in which heaven itself established the shrine.\endnote{Sahag\'un's
appendix is quoted and discussed at length by Garc\'ia Icazbalceta,
\work{Carta}, \S\S17--26 (same witness), which also notes that the
passage stands identically in the editions of Carlos Mar\'ia de
Bustamante and of Lord Kingsborough. Sahag\'un's text was not
collated for this study in the Florentine Codex or in a critical
edition; the passage is reported at the level of the nineteenth-century
critical literature and of its universal citation in the modern
literature on both sides.}

\subsection{September 1556: Mont\'ufar, Bustamante, and the
\work{Informaci\'on}}

The single richest sixteenth-century file bearing on the cult was
generated by a quarrel between the second archbishop of Mexico and
the Franciscan provincial.

On or about 6 September 1556 Archbishop Alonso de Mont\'ufar
preached in his cathedral in support of the devotion at Tepeyac. On
Tuesday 8 September, the feast of the Nativity of Our Lady, at the
solemn Mass in the chapel of San Jos\'e de los Naturales in the
Franciscan convent, with the viceroy Luis de Velasco, the judges of
the \term{audiencia} and a great crowd present, the provincial
Francisco de Bustamante preached. After an unexceptionable panegyric
of the Mother of God he changed subject, visibly agitated, and
attacked the Guadalupe devotion.

The substance of what he said is recoverable because the archbishop
began, the very next day, to take testimony about it. Nine witnesses
deposed. Their evidence, as reconstructed in Edmundo O'Gorman's
close analysis of the file, includes the following: that the
devotion which the people of the city had taken up at ``a hermitage
and house of Our Lady which they have entitled of Guadalupe'' was
gravely harmful to the Indians; that the friars had laboured to teach
the Indians not to adore images, which are wood and stone; that the
archbishop was wrong to think the Indians were not devoted to Our
Lady---they were so devoted that they took her for God; that had the
archbishop understood the Indians as the friars did, ``he would have
taken another manner or order at the beginning of this devotion of
this hermitage''; and, in the sentence that has echoed through four
centuries of controversy, that to tell the Indians now that an image
painted yesterday by an Indian worked miracles was to sow great
confusion and undo what had been well planted.

\begin{witnesstext}{Testimony to the sermon of Fray Francisco de
Bustamante, 8 September 1556, as reconstructed and quoted by Edmundo
O'Gorman}
``y venir ahora a decirles a los naturales que una imagen pintada
ayer por un indio llamado Marcos hac\'ia milagros, era sembrar gran
confusi\'on y deshacer lo bueno que se hab\'ia plantado''\par
``que estaba admirado de que esta devoci\'on se hubiere levantado tan
sin fundamento''\par
``que se maravillaba mucho de que el se\~nor arzobispo hubiese
predicado en los p\'ulpitos y afirmado los milagros que se dec\'ia
que la dicha imagen hab\'ia hecho, siendo prohibido\ldots\ [y que]
era menester haber verificado los milagros y comprob\'adolos con
copia de testigos''\endnote{Edmundo O'Gorman, \work{Destierro de
sombras. Luz en el origen de la imagen y culto de Nuestra Se\~nora de
Guadalupe del Tepeyac}, 2nd ed.\ (Mexico: Universidad Nacional
Aut\'onoma de M\'exico, Instituto de Investigaciones Hist\'oricas,
Serie Historia Novohispana~36), part~II, ch.~3, ``El serm\'on del
provincial fray Francisco de Bustamante,'' pp.~81--92, at
\S\S II.8, II.7 and II.11, read 2026-07-25 in the UNAM open-access
PDF,
\url{https://historicas.unam.mx/publicaciones/publicadigital/libros/222c/222c_04_02_03_CapituloTercero.pdf}.
O'Gorman's own citations are to the \work{Informaci\'on de 1556},
pp.~217, 219, 227, 232, 236, 242--243, 250 (the ``pintada ayer''
sentence), p.~218 (``sin fundamento''), pp.~226 and 242--243 (the
demand for verification). He records that of the witnesses only
Alonso S\'anchez de Cisneros gave the painter's name as Marcos. The
UNAM digital edition carries the notice: ``Se autoriza la
reproducci\'on sin fines lucrativos, siempre y cuando no se mutile o
altere; se debe citar la fuente completa y su direcci\'on
electr\'onica.'' The quoted matter is sixteenth-century testimony
transcribed in a modern critical edition; the surrounding
reconstruction is O'Gorman's and is attributed as such. This study
did not consult the \work{Informaci\'on} in manuscript or in any of
its printed editions.}
\end{witnesstext}

Bustamante went further: he asked the viceroy and the judges to
investigate, to find and flog the ``first inventor'' who had
published that the image worked miracles, and to audit the alms
collected at the hermitage. The archbishop's inquiry stopped without
a sentence; a marginal note in his hand on folio~5 recto reads
``Susp\'endase y la parte es muerto.''\endnote{O'Gorman,
\work{Destierro de sombras}, Appendix~IV, ``La \work{Informaci\'on
de 1556} no es proceso. Sus irregularidades. Su \'indole de documento
no oficial,'' pp.~227--238, at \S\S I and II.14, read 2026-07-25,
\url{https://historicas.unam.mx/publicaciones/publicadigital/libros/222c/222c_04_05_04_ApendiceCuarto.pdf}.
The petitions to the civil power are at part~II, ch.~3, \S\S II.12--14.}

O'Gorman's Appendix~IV establishes two things about this file that
matter for how it may be used, and they cut in different directions.
Against the apparitionist tradition, he shows in detail that the
inquiry was riddled with procedural irregularity by the archbishop's
own diocesan ordinances---anonymous unsigned denunciations, an
interrogatory drawn up not by the fiscal or apostolic notary but by
the viceroy's chaplain, no list of witnesses, oaths largely
administered by the archbishop himself rather than the notary,
leading questions, a witness compelled under pain of excommunication,
another who was the archbishop's own servant, and named persons never
called---so that the flattering picture of a magnanimous prelate who
tried and then forgave the friar cannot stand. Against a careless use
of the file as a simple deposition about Guadalupe, he shows equally
that it is not a process at all but an informative inquiry that the
archbishop assembled for his own protection and kept private; which
is why, as Appendix~VII documents, nobody knew of its existence for
nearly three centuries.\endnote{O'Gorman, Appendix~IV, \S\S I--III;
Appendix~VII, ``Hallazgo y divulgaci\'on de la \work{Informaci\'on de
1556},''
\url{https://historicas.unam.mx/publicaciones/publicadigital/libros/222c/222c_04_05_07_ApendiceSeptimo.pdf},
read 2026-07-25. Appendix~IV also corrects, with the sources, the
claim that Bustamante fell into disgrace for the sermon: he was still
provincial on 6 June 1557, was elected provincial again three years
later, was made commissary general in September 1561, went to court
in 1562 on the Order's business, and died in Madrid on 1 November
1562.}

The transmission history of the \work{Informaci\'on} is itself part
of the evidence and is unusually well documented. On O'Gorman's
reconstruction, drawn from the 1888 letter of Jos\'e Mar\'ia de
Agreda y S\'anchez: Archbishop Manuel Posada y Gardu\~no, shortly
before his death on 30 April 1846, laid his hand on the file and told
the historian Jos\'e Fernando Ram\'irez that what was certain about
the matter was in it, ``but neither you nor anyone else is to see
it,'' and had it placed in the reserved archive. During the ensuing
vacancy the archdeacon Jos\'e Braulio Sagaceta found it, could not
read the old hand well, inferred that it was against the Guadalupan
history, took it home, and kept it hidden for more than twenty years.
It reached Garc\'ia Icazbalceta about 1871. It was first cited in his
letter of 1883, and first printed in 1888, in an edition financed by
opponents of the tradition and falsely imprinted ``Madrid.'' A second
edition of 1890, by Fortino Hip\'olito Vera, printed the same
document under a title asserting that it \emph{proved} the
apparition.\endnote{O'Gorman, Appendix~VII, nos.~1, 3, 11, 13, 20,
36, 38, 42, and 44 (same witness). The 1888 first edition is
\work{Informaci\'on que el arzobispo de M\'exico D.~Fray Alonso de
Mont\'ufar mand\'o practicar\ldots}, imprinted Madrid, Imprenta de la
Guirnalda, but in fact printed in Mexico at the expense of the canon
Vicente de P.~Andrade; the second edition is Vera's \work{La
milagrosa aparici\'on de Nuestra Se\~nora de Guadalupe, comprobada
por una informaci\'on levantada en el siglo XVI} (Amecameca, 1890);
the third is the Irineo Paz edition of 1891.}

That last fact is the fairest possible summary of the state of the
question. The same file was published within three years by both
parties, each as a decisive proof of its own thesis.

\subsection{The rejoinders, stated at their strength}

The argument from silence is not self-evidently decisive, and the
serious apparitionist replies are these.

\emph{The Franciscan reply.} The order that dominated early Mexican
evangelization was, on the evidence of the 1556 file itself, hostile
to a cult of images among neophytes. Silence from Motolin\'ia,
Mendieta and Sahag\'un about a devotion they distrusted is not the
same as silence about an event they had never heard of, and
Bustamante's sermon is evidence of the hostility, not of ignorance.
O'Gorman himself concedes that the file shows the Mexican Franciscans
sharing their provincial's view.\endnote{O'Gorman, Appendix~IV,
\S I, citing the testimonies of Gonzalo de Alarc\'on, Alonso
S\'anchez de Cisneros and Juan de Masseguer, \work{Informaci\'on de
1556}, pp.~239--244, 248--249, and referring to Sahag\'un's known
pronouncement.}

\emph{The transmission reply.} The narrative is Nahuatl, Indigenous,
and oral in the first instance. Alphabetic writing in Nahuatl was a
missionary import; the genres in which such a memory would first be
carried---annals, songs, pictorial records---survive fragmentarily
and were not read by Spanish chroniclers. A story preserved among
Indians in Nahuatl could be unknown to the Spanish record for
decades without ceasing to be true.

\emph{The incompleteness reply.} Garc\'ia Icazbalceta concedes the
premise on which this reply builds: ``we do not know everything that
came from his pen, nor is it reasonable to demand so much.'' His
answer is that if nothing appears in the large surviving body, it is
gratuitous to postulate it in the lost remainder---but that is a
judgment of probability, not a demonstration.\endnote{Garc\'ia
Icazbalceta, \work{Carta}, \S12, same witness.}

\emph{The interested-silence reply.} A conquered people's memory of a
heavenly favour to one of their own was not an obviously safe thing
to publish in the decades of the \work{encomienda} disputes, the
epidemics, and the campaigns against idolatry.

The counter-arguments to each of these are equally available, and
this study does not adjudicate among them. What it does hold is that
the rejoinders change the shape of the problem rather than dissolving
it. The Franciscan reply explains the friars but not Zum\'arraga; the
transmission reply explains the Spanish record but must then account
for the absence of a dated Nahuatl witness; the incompleteness reply
is a possibility rather than a piece of evidence. And the 1556 file
remains what it is: the one occasion on which the origin of the
Tepeyac image was formally at issue in the sixteenth century, in
which a provincial said in a cathedral pulpit before the viceroy that
an Indian had painted it, and in which, on the reading of every
critical historian who has examined the file, no one is recorded as
having contradicted him with the story of Juan Diego.

\subsection{The viceroy's letter of 1575}

One further document is regularly joined to the 1556 file. On 23
September 1575 Viceroy Mart\'in Enr\'iquez wrote to the king an
account of the origin of the Tepeyac devotion in which he reported
that it began when a herdsman published that he had recovered his
health by going to the hermitage. O'Gorman treats the coincidence
between that report and Bustamante's demand that the ``first
inventor'' of the miracle stories be punished as strong confirmation
that the cult's efflorescence belongs to 1555--1556; Garc\'ia
Icazbalceta had drawn the same inference ninety years
earlier.\endnote{O'Gorman, part~II, ch.~3, \S II.13 and
\S III, citing ``Carta del virrey don Mart\'in Enr\'iquez al rey,
M\'exico, 23 de septiembre de 1575,'' in \work{Cartas de Indias},
p.~310; Garc\'ia Icazbalceta, \work{Carta}, \S32. The viceroy's
letter was not consulted at first hand for this study; it is reported
at the level of these two witnesses, which agree on its content.}

Here too the honest statement is bounded. The letter is a viceroy's
report of what he had been told about an origin two decades before
his own arrival; it is excellent evidence that in 1575 the Crown's
chief officer in New Spain believed the cult to be recent and to have
begun with a cure, and it is not a transcript of 1531.
